
\chapter{Hyponymy}
After dealing with various issues of polysemy, some of the second-tier semantic relations may have to be redefined when looked at from a distributional perspective, e.g. hyponymy \parencite[443]{herbelotMeasuringSemanticContent2013}. In this view, other existing taxonomies that may have acted as models of hyponymy need to be replaced by a definition that reflects how actual usage contexts can limit our selection of possible candidates for hyponymy \parencite[443]{rollerInclusiveSelectiveSupervised2014}. Unfortunately, this view may clash with the observations of \textcite{herbelotMeasuringSemanticContent2013}:
\begin{itemize}
      \item Hypernyms can have very specific usage contexts.
      \item Collocations may lead to unexpected results in the distribution of words.
\end{itemize}
Similar adjustments have to be made in treating other kinds of semantic relations. Therefore, many researchers apply preprocessing to their corpora before performing the actual analysis on them \parencite[56f.]{gyllenstenRgramsUnsupervisedLearning2019}. However, this increases the danger of circularity, revealing only those results that we were already expecting in the first place. As an alternative, we should look for a model that is robust enough to detect semantic relations despite the confounding factors. The objective should be to fit the model to the base text, not the other way round. For vector space models like Skip-gram \parencite[5]{mikolovEfficientEstimationWord2013}, subsampling and the restriction of strong collocations may be justified from a computational perspective because they enable a faster training of models while more or less preserving the resulting performance in the prediction of word contexts. This very tradeoff, though, may not work out as well for the detection of semantic relations, e.g. if we give a higher priority to the frequency of certain contexts. Such a change of algorithm may particularly be motivated by the distinction of theoretical and empirical scopes of words \parencite[46]{coenenAnalogieUndMetapher2013}. 

\label{distributionalInclusionHypothesis}
This also applies to the domains of corpus linguistics and language learning: Given a text corpus that helps learners to acquire a certain knowledge (e.g. all the ancient literature mentioned in an educational curriculum for the Classics), the theoretical scope of a certain word would be equivalent to the set of all distinct contexts in the corpus in which that word occurs. In contrast, the empirical scope of that same word would be equal to a subset of that theoretical scope, namely those contexts which are already known to the learner (at any given point in time of the learning process). Thus, the theoretical scope will remain static as long as we do not change the corpus, while the empirical one will be dynamic as long as we keep learning. From this point of view, we can also restrict our notion of semantic relations: Hyponymy, for instance, will be tied to those lexical networks that are present in a specific corpus and, for a specific learner, specific instances of hyponomy relations will not exist unless they have already been encountered in the corpus and lead to a relevant learning process. It is therefore our duty to a) detect such instances in the corpus and b) use them to design learning situations that help language learners to construct lexical and semantic knowledge \parencite[25]{yorkConstructivistApproachGamebased2018}. In order to accomplish that first part, we need to have clearly defined requirements to qualify a relation as one of hyponymy, e.g. according to a version of the Distributional Inclusion Hypothesis \parencite[443]{rollerInclusiveSelectiveSupervised2014} that also accounts for idiomaticity, burstiness and other kinds of lexical interferences.

\section{WordNet}
\label{WordNet}
Lexical resources like the Latin WordNet \parencite{minozziLatinWordNetProject2010} may contribute a reference model for such a purpose because it includes a built-in notion of hyponymy. Unfortunately, the WordNet defines such relations based on intuition, not on actual language use \parencite[315]{fellbaumChallengesMultilingualWordnet2012}. WordNet is therefore static and non-empirical from the perspective of usage-based fields like corpus linguistics. However, it was constructed by alleged experts and we may use their assessment for reflecting our own approach. This kind of triangulation is supported by recent studies, which often explicitly combine knowledge bases and corpus data \parencite[984-988]{onoWordEmbeddingbasedAntonym2015}, thereby integrating both the introspective and the usage-based approaches to semantics. The former can be helpful in two ways: a) by defining a horizon of expectations in language learning contexts, i.e. which information learners are supposed to infer from their previous knowledge in order to understand a linguistic expression in a specific text passage; and b) by modeling semantic knowledge in a machine-readable, consistent (albeit simplistic) manner. The integration of distributional semantics, on the other hand, is particularly motivated by the hypothesis that semantic properties of words are reflected in their contextual surroundings \parencite[59]{griesBehavioralProfilesCorpusbased2009}. 

\label{transitivity}
By using distributional data from corpora, we also gain the advantages of rigorously empirical research (as opposed to bare intuition and anecdotal evidence): frequency, authenticity, variation and systematic induction \parencite[60]{griesBehavioralProfilesCorpusbased2009}. Still, knowledge databases may be useful for languages like Latin where previous research is scarce and, more often than not, heavily intuition-based or anecdotal. In both the introspective and the distributional case, we can increase the amount of provided information by relying on associated properties of hyponymy relations, e.g. transitivity \parencite[58]{coenenAnalogieUndMetapher2013}. Applied to the Latin WordNet, this may look as follows:
\begin{enumerate}
         \item \textit{gladius}
         \item a cutting or thrusting weapon with a long blade
         \item weaponry used in fighting or hunting
         \item weapons considered collectively
         \item an artifact (or system of artifacts) that is instrumental in accomplishing some end
         \item a man-made object
         \item a physical (tangible and visible) entity
\end{enumerate}
This information flow can be read from both sides. Beginning at \textit{gladius} (sword), it can be used to simulate the process of reading a text: We encounter a word in the text (e.g. the plural \textit{gladii}) and determine its lemma (\textit{gladius}). From there, we can try to find a corresponding meaning for the surface form, e.g. "a cutting or thrusting weapon with a long blade". This meaning, in the WordNet taxonomy, can be hyponym to other kinds of meaning, e.g. "weaponry used in fighting or hunting". Now, if we know that
\begin{itemize}
         \item \textit{gladius} is a weapon with a blade and
         \item weapons with a blade can be used for hunting,
\end{itemize}
then we are safe to assume that a \textit{gladius} can be used for hunting. The further we go down the list, the more abstract do the meanings get. Nonetheless (or because of that), every layer can remind us of new specific features that we may want to associate with \textit{gladius}: specific motion (2.), specific purpose (3.), hypernym (4.), general purpose (5.), origin and production (6.), material (7.). In some cases, we may even use morphological hints to get faster access to information about hyponymy \parencite[913]{anstattTypenSemantischerRelationen2009}. This is especially true for Latin with its relatively rich morphology: the meaning of \textit{artifex} (artist) may be associated directly to \textit{ars} (art) and \textit{facere} (to make), without necessarily traversing a hierarchical tree-like representation of increasing semantic generalization. Interestingly, this is possible even in cases where the compound is formed from foreign language material and thus cannot be easily paraphrased in the native language \parencite[33]{souille-rigautSemanticAccountQuasiLexemes2010}. 

How is that important for Latin hyponymy? Consider Greek loanwords: \textit{xylophytum} (a certain kind of plant) cannot be reasonably described in Latin words as a \textit{*phytum} that is made of \textit{*xylum}. However, Romans may have deduced that the meaning of \textit{xylophytum} is related to ($\rightarrow$ hyponymy) \textit{-phyt-}, thereby inferring previous knowledge about plants. Therefore, we have to allow for the possibility that some instances of hyponymy may not be extractable solely on a distributional basis, but rather on a morphological one. Besides, a morphologically informed model is able to explain why we can immediately see semantic structure in newly coined words \parencite[43]{souille-rigautSemanticAccountQuasiLexemes2010}. In this respect, knowledge about specific hyponyms or hyponymy in general can facilitate the language learning process by providing a rather intuitive access to the meaning of single words or phrases. Unfortunately, this intuition sometimes competes with other interpretations of a word's sense which may be equally intuitive and valid \parencite[126]{ponsborderiaPathsGrammaticalizationSpanish2014}. 

Such examples underline the importance of disambiguation (\hyperref[polysemyProblemForSemanticRelations]{see above}), which may not always be possible. Furthermore, additional steps may be necessary even if a specific word sense was already determined, e.g. the inference of common knowledge or common sense \parencite[427f.]{pinkalSemantik1993}. In this view, disambiguation should be preceded by a thorough linguistic analysis of the context, which, in terms of distributional semantics, is usually addressed by extracting cooccurrence frequencies from sliding windows over a textual input, be it plain text or a linguistically annotated treebank. However, the integration of external information into this process is not as straightforward and therefore often avoided. In the end, this might not be the worst decision, depending on our definition of semantics: Some models of distributional semantics reduce the spectrum of polysemy to one main aspect of meaning per word.
\parencite[4]{faruquiProblemsEvaluationWord2016}.

\hypertarget{ELMO}{Approaches} of this kind level nuances that may actually be important for our understanding of a single word in its context. As an improvement, other models regard entire sentences for the representation of a word's context-specific meaning, thereby differentiating between various common usages \parencite[2f.]{petersDeepContextualizedWord2018}. Modeling the process of language production both in a forward and backward manner at the same time seems reasonable, especially for the Latin language, because the antecedent of, e.g., a relative pronoun can be part of the relative clause, which makes it difficult to understand the structure of such a construction by reading it in just one direction, without referring to the past and future context at the same time:\footnote{The glossing of longer Latin text passages was created according to the \href{https://www.eva.mpg.de/lingua/pdf/Glossing-Rules.pdf}{Leipzig Glossing Rules}.}

\vspace{10pt}
\ex[exno=PROIEL 53467]
	\par
	\begingl
		\gla nam et frument-um ex agr-is cotidie in castr-a confere-ba-t et qu-ae grav-issim-e adfli-ct-ae era-nt nav-es e-a-rum materi-a atque ae-re ad reliqu-a-s reficie-nd-a-s ute-bat-ur [...] //
		\glb since also grain-{\sc acc} from field-{\sc abl.pl} daily-{\sc adv} into camp-{\sc acc.pl} convey-{\sc impf-3sg} and which-{\sc f.pl} heavy-{\sc sup-adv} strike-{\sc pass-f.pl} be-{\sc 3pl} ship-{\sc pl} this-{\sc f-pl.gen} timber-{\sc abl} and brass-{\sc abl} to remaining-{\sc f-pl.acc} repair-{\sc gdv-f-pl.acc} use-{\sc impf-3sg}//
		\glft for he daily conveyed corn from the country parts into the camp, used the timber and brass of such ships as were most seriously damaged for repairing the rest [...] \parencite{mcdevitteCaesarGallicWar1869} //
	\endgl
\xe
\vspace{10pt}

In this example, \textit{naves} is the antecedent to the relative pronoun \textit{quae}. Its integration into the relative clause is indicated by its morphological congruency (case, number, gender) with the relative pronoun. In a textbook setting, the sentence would rather look like this: \textit{[...] conferebat et earum navium, quae gravissime adflictae erant, materia atque [...]}. Such syntactic variations are probably related to the relatively free word order in the Latin language. Therefore, a language model for Latin should be flexible enough to analyze a given textual sequence by looking at either the left or right context, according to the circumstances of the case \parencite[8]{devlinBERTPretrainingDeep2019}. While this strictly bidirectional, attention-based approach \parencite{vaswaniAttentionAllYou2017} enhances the ability to explain word meanings from their larger surrounding contexts, it does not address the problem of inference, e.g. relevant external text passages. To analyze this problem empirically, we will set up a toy language model in order to inspect difficult cases.